%% toc: P3 34 34.9
\section{Proofs for real closed fields as the principal worked example}
\label{sec:annex-p03-w03-p03-037-rcf}

Throughout, \(K\) is a real closed field in the sense of
\cref{def:p03-w03-p03-037-rcf}.  The ordered-field facts used without
comment are: \(0<1\) (if \(1<0\) then \(0<-1\) and \(0<(-1)(-1)=1\)); squares
are nonnegative; \(n\cdot1>0\) for every integer \(n\geq1\); if
\(\abs{E}<\abs{w}\) then \(w+E\) has the sign of \(w\); and products of
nonzero elements are nonzero.  The intermediate-value axiom of degree \(n\)
quantifies over all coefficient tuples, so it applies to every polynomial
of degree at most \(n\) that changes sign.

\begin{annexitem}{annex:p03-w03-p03-037-poly-basics}{thm:p03-w03-p03-037-poly-basics}{One-variable polynomial facts}
\begin{proof}
(1) By induction on \(p\).  For \(p=0\), \(P=a_0\neq0\) has no root.  If
\(p\geq1\) and \(c\) is a root, division by the monic polynomial \(X-c\)
gives \(P=(X-c)Q+P(c)=(X-c)Q\) with \(Q\) of degree \(p-1\).  A root
\(d\neq c\) of \(P\) satisfies \((d-c)Q(d)=0\), hence \(Q(d)=0\).  By the
induction hypothesis \(P\) has at most \(1+(p-1)\) roots.

(2) For \(p=0\) the sign is constant.  Let \(p\geq1\) and suppose
\(u<v\) in \(C\) with \(\operatorname{sgn}P(u)\neq\operatorname{sgn}P(v)\).
Neither value is \(0\), since \(C\) contains no root, so
\(P(u)P(v)<0\).  The intermediate-value axiom gives a root in \((u,v)\), and
\((u,v)\subseteq C\) because \(C\) is convex; a contradiction.

(3) If \(p=0\) then \(M_P=1\) and \(\operatorname{sgn}P(x)=\operatorname{sgn}a_0\).
Let \(p\geq1\) and \(\abs x>M_P\geq1\).  Then \(\abs x^k\leq\abs x^{p-1}\)
for \(k\leq p-1\), so
\[
\Bigl\lvert\sum_{k<p}a_kx^k\Bigr\rvert\leq\sum_{k<p}\abs{a_k}\abs x^{p-1}
=(M_P-1)\abs{a_p}\abs x^{p-1}<\abs x\abs{a_p}\abs x^{p-1}=\abs{a_px^p}.
\]
Hence \(P(x)=a_px^p+E\) with \(\abs E<\abs{a_px^p}\), and
\(\operatorname{sgn}P(x)=\operatorname{sgn}(a_px^p)=
\operatorname{sgn}(a_p)\operatorname{sgn}(x)^p\).

(4) The roots of \(P\) in \((a,b]\) form a nonempty set (it contains \(b\))
that is finite by (1); let \(d\) be its least element and \(c=a\).  Then
\(P(c)=P(d)=0\) and \(P\) has no root in \((c,d)\).  Dividing out the
largest powers of \(X-c\) and \(X-d\) as in (1),
\(P=(X-c)^m(X-d)^nQ\) with \(m,n\geq1\) and \(Q(c)\neq0\neq Q(d)\).  Every
root of \(Q\) is a root of \(P\), so \(Q\) has no root in the convex set
\([c,d]\), and by (2) \(\operatorname{sgn}Q(c)=\operatorname{sgn}Q(d)=s\neq0\).
By the product rule for formal derivatives,
\(P'=(X-c)^{m-1}(X-d)^{n-1}R\) with
\[
R=m(X-d)Q+n(X-c)Q+(X-c)(X-d)Q'.
\]
Then \(R(c)=m(c-d)Q(c)\) has sign \(-s\) and \(R(d)=n(d-c)Q(d)\) has sign
\(s\), because \(m\cdot1,n\cdot1>0\) and \(c<d\).  So \(R(c)R(d)<0\), and the
intermediate-value axiom gives \(e\in(c,d)\) with \(R(e)=0\); then
\(P'(e)=0\) and \(e\in(a,b)\).

(5) Put \(\lambda=(P(b)-P(a))/(b-a)\) and
\(G=P-P(a)-\lambda(X-a)\), so \(G(a)=G(b)=0\).  If \(G\) is the zero
polynomial, then \(P'=\lambda\) identically and \(c=(a+b)/2\) works.
Otherwise (4) applied to \(G\) gives \(c\in(a,b)\) with
\(0=G'(c)=P'(c)-\lambda\).

(6) The coefficient of \(X^{p-1}\) in \(P'\) is \(p\,a_p\neq0\), so \(P'\) is
not the zero polynomial, and by (2) it has a constant sign \(\varepsilon\)
on \(J\), which is nonzero because \(P'\) has no root in \(J\).  If
\(x<x'\) lie in \(J\) or are endpoints of \(J\) in \(K\), then
\((x,x')\subseteq J\), and (5) gives \(c\in(x,x')\) with
\(P(x')-P(x)=(x'-x)P'(c)\), whose sign is \(\varepsilon\).

(7) Let \(c\in K\) and \(\abs{x-c}\leq1\).  Since
\(x^k-c^k=(x-c)\sum_{i<k}x^ic^{k-1-i}\) and \(\abs x\leq\abs c+1\),
\(\abs{x^k-c^k}\leq k(\abs c+1)^{k-1}\abs{x-c}\).  Hence
\(\abs{P(x)-P(c)}\leq B\abs{x-c}\) with
\(B=\sum_{k\leq p}k\abs{a_k}(\abs c+1)^{k-1}\).  Given an open interval
around \(P(c)\), choose \(\eta>0\) such that \((P(c)-\eta,P(c)+\eta)\) lies
inside it; for \(\abs{x-c}<\min(1,\eta/(B+1))\) we get
\(\abs{P(x)-P(c)}<\eta\).  Since open intervals form a basis of the order
topology, \(P\) is continuous.

For square roots, let \(x\geq0\).  If \(x=0\) take \(0\).  If \(x>0\), the
polynomial \(Y^2-x\) is \(-x<0\) at \(0\) and \(x^2+x+1>0\) at \(x+1\), so the
intermediate-value axiom gives \(y\in(0,x+1)\) with \(y^2=x\).  If
\(y,y'\geq0\) and \(y^2=y'^2\), then \((y-y')(y+y')=0\); if \(y+y'=0\) both
are \(0\), so in either case \(y=y'\).  The map \(\sqrt{\cdot}\) is strictly
increasing on \([0,\infty)\): if \(0\leq r'<r\) then \(r'^2<r^2\).  For
\(u,v\geq0\) let \(r=\sqrt u\) and \(r'=\sqrt v\), say \(r\geq r'\).  Then
\((r-r')^2\leq(r-r')(r+r')=u-v=\abs{u-v}\), and since both \(r-r'\) and
\(\sqrt{\abs{u-v}}\) are nonnegative, \(r-r'\leq\sqrt{\abs{u-v}}\).  Given
\(\eta>0\), \(\abs{u-v}<\eta^2\) gives
\(\abs{\sqrt u-\sqrt v}\leq\sqrt{\abs{u-v}}<\eta\), so \(\sqrt{\cdot}\) is
continuous.
\end{proof}
\end{annexitem}

\begin{annexitem}{annex:p03-w03-p03-037-interval}{thm:p03-w03-p03-037-interval}{Behaviour of a polynomial between roots of its derivative}
\begin{proof}
By \cref{thm:p03-w03-p03-037-poly-basics}(6), \(\varepsilon\neq0\) and
\(\varepsilon(P(x')-P(x))>0\) whenever \(x<x'\) belong to
\(\bar J=J\cup(\{\alpha,\beta\}\cap K)\).  Since \(a_p\neq0\),
\(s_\alpha\neq0\) when \(\alpha=-\infty\) and \(s_\beta\neq0\) when
\(\beta=+\infty\).

\emph{Witness points.}  For \(x\in J\) let \(u_x=\alpha\) if
\(\alpha\in K\) and \(u_x=\min(x,-M_P)-1\) if \(\alpha=-\infty\); let
\(v_x=\beta\) if \(\beta\in K\) and \(v_x=\max(x,M_P)+1\) if
\(\beta=+\infty\).  Then \(u_x<x<v_x\), \(u_x,v_x\in\bar J\), and by
\cref{thm:p03-w03-p03-037-poly-basics}(3)
\(\operatorname{sgn}P(u_x)=s_\alpha\) and \(\operatorname{sgn}P(v_x)=s_\beta\).
Hence
\[
\varepsilon\bigl(P(x)-P(u_x)\bigr)>0,\qquad
\varepsilon\bigl(P(v_x)-P(x)\bigr)>0.\tag{\(*\)}
\]

\emph{The cases are exclusive and exhaustive.}  If \(s_\alpha=s_\beta=0\),
both endpoints lie in \(K\) and \(P(\alpha)=P(\beta)=0\), contradicting
\(\varepsilon(P(\beta)-P(\alpha))>0\).  So at most one of
\(s_\alpha,s_\beta\) is \(0\); if neither is, they are equal or opposite.

(1) If \(s_\alpha=0\), then \(\alpha\in K\), \(u_x=\alpha\) and
\(P(u_x)=0\), so \((*)\) gives \(\operatorname{sgn}P(x)=\varepsilon\) for
every \(x\in J\).

(2) If \(s_\beta=0\), symmetrically \(\operatorname{sgn}P(x)=-\varepsilon\).

(3) If \(s_\alpha=s_\beta=s\neq0\), then by \((*)\) the value \(P(x)\) lies
strictly between \(P(u_x)\) and \(P(v_x)\), both of sign \(s\); so
\(P(x)\) has sign \(s\).  In particular \(P\) has no root in \(J\).

(4) If \(s_\alpha=-s_\beta\neq0\), fix \(x_0\in J\) and write \(u=u_{x_0}\),
\(v=v_{x_0}\).  Then \(P(u)P(v)<0\), and the intermediate-value axiom gives
\(c\in(u,v)\subseteq J\) with \(P(c)=0\).  By strict monotonicity on
\(\bar J\), \(c\) is the only root in \(J\).  For \(x\in(\alpha,c)\),
\(\varepsilon(P(c)-P(x))>0\) gives \(\operatorname{sgn}P(x)=-\varepsilon\);
and since \(u<c\), \(\varepsilon(P(c)-P(u))>0\) gives
\(s_\alpha=\operatorname{sgn}P(u)=-\varepsilon\).  So \(P\) has sign
\(s_\alpha\) on \((\alpha,c)\).  Symmetrically, for \(x\in(c,\beta)\),
\(\operatorname{sgn}P(x)=\varepsilon=\operatorname{sgn}P(v)=s_\beta\).
\end{proof}
\end{annexitem}

\begin{annexitem}{annex:p03-w03-p03-037-sublist}{thm:p03-w03-p03-037-sublist}{Sign tables of sublists}
\begin{proof}
For a sign vector \(\sigma\in\{-1,0,1\}^s\) let \(\pi\sigma\) be the vector
\((\sigma_{\iota(1)},\dots,\sigma_{\iota(r)})\).  Given a word
\(w=(w_0,w_1,\dots,w_{2N})\) of odd length, let
\(E=\{i\leq r:(w_0)_{\iota(i)}\neq0\}\), and let \(j_1<\dots<j_M\) be the
indices \(j\in\{1,\dots,N\}\) for which \((w_{2j-1})_{\iota(i)}=0\) for some
\(i\in E\).  Define
\[
\rho_\iota(w)=\bigl(\pi w_0,\ \pi w_{2j_1-1},\ \pi w_{2j_1},\ \dots,\
\pi w_{2j_M-1},\ \pi w_{2j_M}\bigr),
\]
which is \((\pi w_0)\) when \(M=0\), and define \(\rho_\iota\) arbitrarily on
words of even length.

Now let \(w=\operatorname{SIGN}(\mathcal P)\), with roots
\(z_1<\dots<z_N\).  A member \(P_k\) is the zero polynomial exactly when
\((w_0)_k=0\): the cell \(I_0\) is infinite (it is \(K\), or contains
\(z_1-n\cdot1\) for all \(n\geq1\)), so a nonzero \(P_k\) has a non-root in
\(I_0\), and its constant sign there is nonzero.  Hence \(E\) indexes the
nonzero members of the sublist.  Every root of a nonzero member of the
sublist is a root of a nonzero member of \(\mathcal P\), hence some
\(z_j\); and \(z_j\) is such a root exactly when \(j\in\{j_1,\dots,j_M\}\).
So the roots of the sublist are \(z_{j_1}<\dots<z_{j_M}\).  Its cells are
\(\{z_{j_k}\}\), whose sign vector is \(\pi w_{2j_k-1}\), and the intervals
\((-\infty,z_{j_1})\), \((z_{j_k},z_{j_{k+1}})\) and
\((z_{j_M},+\infty)\), which contain the nonempty cells \(I_0\),
\(I_{j_k}\) and \(I_{j_M}\) of \(\mathcal P\) respectively.  The sublist has
constant signs on each of its cells, so its sign vector there equals the
projection of the sign vector on the contained cell of \(\mathcal P\), namely
\(\pi w_0\), \(\pi w_{2j_k}\) and \(\pi w_{2j_M}\).  This is
\(\rho_\iota(w)\).  When \(\iota\) is a bijection, every root is kept and
\(\rho_\iota\) only permutes coordinates.
\end{proof}
\end{annexitem}

\begin{annexitem}{annex:p03-w03-p03-037-prem}{thm:p03-w03-p03-037-prem}{Pseudo-division}
\begin{proof}
By induction on \(p\geq q-1\).  If \(p=q-1\), take \(A=0\) and \(R=P\).  If
\(p\geq q\), write \(P=a_pX^p+P_{\leq p-1}\) and put
\[
P_1=b_q\,P_{\leq p-1}-a_pX^{p-q}\,Q_{\leq q-1}=b_qP-a_pX^{p-q}Q,
\]
which has formal degree \(p-1\geq q-1\) because the terms in \(X^p\) cancel.
By the induction hypothesis \(b_q^{\,p-q}P_1=A_1Q+R_1\) with \(R_1\) of
formal degree \(q-1\).  Then
\[
b_q^{\,p-q+1}P=b_q^{\,p-q}\bigl(P_1+a_pX^{p-q}Q\bigr)
=\bigl(A_1+b_q^{\,p-q}a_pX^{p-q}\bigr)Q+R_1,
\]
so \(R=R_1\) works.  This recursion fixes \(\operatorname{prem}(P,Q)=R\) as a
definite parametric polynomial.

For the consequence, evaluate the identity at \(\bar y\) and \(w\):
\[
b_q(\bar y)^{p-q+1}P(\bar y)(w)=A(\bar y)(w)\,Q(\bar y)(w)+R(\bar y)(w)
=R(\bar y)(w).
\]  Taking signs and using
\(\operatorname{sgn}(b_q(\bar y))^{2(p-q+1)}=1\) gives the stated formula.
\end{proof}
\end{annexitem}

\begin{annexitem}{annex:p03-w03-p03-037-uniform-sign}{thm:p03-w03-p03-037-uniform-sign}{Uniform sign tables}
\begin{proof}
For a list \(\mathcal P\) let \(\mu(\mathcal P)=(p,c)\), where \(p\) is the
largest formal degree of a member and \(c\) the number of members of that
formal degree (\(\mu=(0,0)\) for the empty list).  We argue by induction on
\(\mu\) in the lexicographic order of \(\mathbb N\times\mathbb N\), which is
a well-order.  Throughout, \(\theta_{\mathcal P,t}\) is \(0\neq0\) for words
\(t\) outside \(\mathbb T_{s,D}\), since such words never occur.

\emph{Base: \(p=0\).}  Every member is a constant \(a_i\in\mathbb Z[\bar Y]\),
so no member has a root and \(\operatorname{SIGN}(\mathcal P(\bar y))\) is
the one-letter word \((\sigma)\) with
\(\sigma_i=\operatorname{sgn}a_i(\bar y)\).  Let \(\theta_{\mathcal P,(\sigma)}\)
be the conjunction of the atomic formulas \(a_i>0\), \(a_i=0\) or
\(a_i<0\) according as \(\sigma_i\) is \(1\), \(0\) or \(-1\), and let
\(\theta_{\mathcal P,t}\) be \(0\neq0\) for every other word.

\emph{Step: \(p\geq1\).}  By \cref{thm:p03-w03-p03-037-sublist} with a
bijection \(\iota\), reordering a list permutes the coordinates of its sign
table by a fixed rule; so we may assume that the last member
\(P=\sum_{k\leq p}a_kX^k\) has formal degree \(p\).  Write
\(\mathcal P=(\mathcal P_0,P)\), let \(\mathcal P^\flat=(\mathcal P_0,P_{\leq p-1})\),
\(\mathcal R=(\mathcal P_0,P')\) and \(\mathcal Q\) as in
\cref{def:p03-w03-p03-037-elimination-step}.  Every member of
\(\mathcal Q\) other than those of \(\mathcal P_0\) has formal degree at most
\(p-1\): \(P'\) has formal degree \(p-1\), the coefficients have formal degree
\(0\), and \(\operatorname{prem}(P,Q_{\leq j})\) has formal degree \(j-1\leq p-1\)
(note \(j\leq q\leq p\), so \cref{thm:p03-w03-p03-037-prem} applies).
\(\mathcal P_0\) has \(c-1\) members of formal degree \(p\).  Hence
\(\mu(\mathcal Q)<\mu(\mathcal P)\), and likewise
\(\mu(\mathcal P^\flat)<\mu(\mathcal P)\).

\emph{Claim.}  There is a function \(\Phi\) on words, depending only on the
positions of the members of \(\mathcal Q\) and \(\mathcal P\), such that
\(\operatorname{SIGN}(\mathcal P(\bar y))=\Phi(\operatorname{SIGN}(\mathcal Q(\bar y)))\)
for every real closed field \(K\) and every \(\bar y\in K^m\) with
\(a_p(\bar y)\neq0\).

Granting the claim, put
\[
\theta_{\mathcal P,t}=
\bigl(a_p=0\wedge\theta_{\mathcal P^\flat,t}\bigr)\vee
\Bigl(\neg\,a_p=0\wedge\bigvee_{u\in\mathbb T_{\mathcal Q},\ \Phi(u)=t}
\theta_{\mathcal Q,u}\Bigr),
\]
where \(\mathbb T_{\mathcal Q}\) is the finite set of words that can be sign
tables of \(\mathcal Q\) (\cref{def:p03-w03-p03-037-sign-table}) and the
\(\theta\) on the right exist by the induction hypothesis.  If
\(a_p(\bar y)=0\) then \(P(\bar y)=P_{\leq p-1}(\bar y)\) in \(K[X]\), so
\(\mathcal P(\bar y)\) and \(\mathcal P^\flat(\bar y)\) are the same list and
have the same sign table.  If \(a_p(\bar y)\neq0\), then
\(\operatorname{SIGN}(\mathcal P(\bar y))=t\) iff
\(\operatorname{SIGN}(\mathcal Q(\bar y))\) is some \(u\) with \(\Phi(u)=t\).
In both cases \(K\entails\theta_{\mathcal P,t}(\bar y)\) iff
\(\operatorname{SIGN}(\mathcal P(\bar y))=t\).

\emph{Proof of the claim.}  Fix \(K\) and \(\bar y\) with
\(a_p(\bar y)\neq0\), and write \(\hat F=F(\bar y)\) for each parametric
polynomial \(F\).  Then \(\hat P\) has degree \(p\geq1\), and \(\hat P'\) has
leading coefficient \(p\,a_p(\bar y)\neq0\), so it is not zero.  Let
\(u=\operatorname{SIGN}(\mathcal Q(\bar y))\); every step below reads only
\(u\) and the fixed positions.

(a) \emph{Roots of \(\mathcal R\).}  Since \(\mathcal R\) is a sublist of
\(\mathcal Q\), \(v=\rho(u)\) is the sign table of \(\hat{\mathcal R}\)
(\cref{thm:p03-w03-p03-037-sublist}); let \(w_1<\dots<w_M\) be its roots and
\(J_0,\dots,J_M\) its open cells, with \(w_0=-\infty\) and
\(w_{M+1}=+\infty\).  By the construction of \(\rho\), the roots of
\(v\) are the kept roots of \(u\) in order, so the letter of \(u\) at each
\(w_k\) is known; and the sign of a constant member of \(\mathcal Q\) is the
same in every letter.

(b) \emph{The sign of \(P\) at each \(w_k\).}  Let \(Q\) be the first member
of \(\mathcal R\) with \((v_0)_Q\neq0\) and \((v_{2k-1})_Q=0\), that is, the
first nonzero member of \(\hat{\mathcal R}\) vanishing at \(w_k\).  Let
\(b_0,\dots,b_q\) be its coefficients, members of \(\mathcal Q\), and let
\(j\) be the largest index with \(\operatorname{sgn}b_j(\bar y)\neq0\).  It
exists because \(\hat Q\neq0\), and \(j\geq1\) because a nonzero constant
has no root.  Then \(\hat Q=\widehat{Q_{\leq j}}\) has leading coefficient
\(b_j(\bar y)\neq0\), and by \cref{thm:p03-w03-p03-037-prem} applied to
\(P\) and \(Q_{\leq j}\),
\[
\tau_k:=\operatorname{sgn}\hat P(w_k)=
\operatorname{sgn}(b_j(\bar y))^{p-j+1}\,
\operatorname{sgn}\widehat{\operatorname{prem}(P,Q_{\leq j})}(w_k),
\]
where the last sign is the coordinate of that pseudo-remainder in the
letter of \(u\) at \(w_k\).

(c) \emph{Signs at infinity.}  Put \(\tau_0=(-1)^p\operatorname{sgn}a_p(\bar y)\)
and \(\tau_{M+1}=\operatorname{sgn}a_p(\bar y)\), read from the coordinate of
the constant \(a_p\).

(d) \emph{Each interval.}  On \(J_k=(w_k,w_{k+1})\), \(\hat P'\) has no root,
since its roots are among the \(w\)'s; let \(\varepsilon_k\) be its
coordinate in \(v_{2k}\), nonzero by
\cref{thm:p03-w03-p03-037-poly-basics}(6).  Apply
\cref{thm:p03-w03-p03-037-interval} to \(\hat P\) on \(J_k\) with
\(s_\alpha=\tau_k\) and \(s_\beta=\tau_{k+1}\): these are the quantities named
there, by (b) and (c).  Exactly one case holds and it is decided by
\(\tau_k,\tau_{k+1}\).  In cases 1--3, \(\hat P\) has no root in \(J_k\) and
has the constant sign \(\varepsilon_k\), \(-\varepsilon_k\) or \(\tau_k\);
in case 4, \(J_k\) splits as \((w_k,c)\), \(\{c\}\), \((c,w_{k+1})\) with
\(\hat P\)-signs \(\tau_k\), \(0\), \(\tau_{k+1}\).  The members of
\(\hat{\mathcal R}\) keep on these pieces the signs \(v_{2k}\) they have on
all of \(J_k\).

(e) \emph{Assembly.}  Form the word whose letters are, in order: for each
\(J_k\) the one or three letters of (d), each obtained by extending
\(v_{2k}\) by the \(\hat P\)-sign of the piece; and for each
\(w_k\) the letter \(v_{2k-1}\) extended by \(\tau_k\).  Its points are the
\(w_k\), which are roots of nonzero members of \(\hat{\mathcal R}\), and the
roots \(c\) of case 4, which are roots of \(\hat P\neq0\).  Conversely every
root of \(\hat P\) is some \(w_k\) or lies in some \(J_k\), where by (d) it
is the root of case 4.  So the word is the sign table of the list
\((\hat{\mathcal P}_0,\hat P',\hat P)\).  Finally let \(\Phi(u)\) be its image
under the map \(\rho\) of \cref{thm:p03-w03-p03-037-sublist} that deletes the
coordinate of \(P'\); that image is \(\operatorname{SIGN}(\hat{\mathcal P})\).
For words \(u\) on which some step is undefined, let \(\Phi(u)=u\); such
words are never sign tables of \(\mathcal Q\) at parameters with
\(a_p(\bar y)\neq0\).  This proves the claim and the theorem.
\end{proof}
\end{annexitem}

\begin{annexitem}{annex:p03-w03-p03-037-qe}{thm:p03-w03-p03-037-qe}{Quantifier elimination for real closed fields}
\begin{proof}
\emph{Atomic formulas.}  By induction on terms, every
\(\Lang_{\mathrm{or}}\)-term denotes in every commutative ring the function of
an integer polynomial, and every integer polynomial is so denoted
(\cref{def:p03-w03-p03-037-parametric}).  In every ordered field
\(t_1=t_2\) is equivalent to \(t_1-t_2=0\) and \(t_1<t_2\) to
\(0<t_2-t_1\).  So every atomic formula is equivalent modulo
\(\mathsf{RCF}\) to a sign condition \(P=0\) or \(P>0\) with \(P\) an integer
polynomial.

\emph{Reduction to one existential quantifier.}  Let \(\mathcal E\) be the
class of formulas that are \(\mathsf{RCF}\)-equivalent to a quantifier-free
formula with the same free variables.  It contains the atomic formulas and
is closed under the Boolean connectives.  Since \(\forall x\,\varphi\) is
equivalent to \(\neg\exists x\,\neg\varphi\), and since
\(\varphi\leftrightarrow\psi\) in all models implies
\(\exists x\,\varphi\leftrightarrow\exists x\,\psi\) in all models, it is
enough to show that \(\exists x\,\psi\in\mathcal E\) for every
quantifier-free \(\psi\); induction on formulas then gives the theorem.

\emph{Normal form.}  Write \(\psi\) in disjunctive normal form.  By the
trichotomy of a linear order, \(\neg(P=0)\) is equivalent to
\(P>0\vee P<0\) and \(\neg(P>0)\) to \(P=0\vee P<0\), where \(P<0\) is an
atomic formula.  Redistributing, \(\psi\) is equivalent modulo
\(\mathsf{RCF}\) to a finite disjunction of conjunctions
\(\bigwedge_{i\leq s}\operatorname{sgn}P_i=\sigma_i\) with
\(\sigma\in\{-1,0,1\}^s\).  The existential quantifier distributes over the
disjunction.

\emph{One conjunction.}  Regard each \(P_i\) as a parametric polynomial in
\(x\), with the remaining free variables \(\bar Y\) as parameters and its
\(x\)-degree as formal degree, and let \(\mathcal P=(P_1,\dots,P_s)\) with
formal degrees summing to \(D\).  For a real closed field \(K\) and
\(\bar y\in K^m\), the cells of \(\mathcal P(\bar y)\) are nonempty, partition
\(K\), and carry constant sign vectors.  Hence
\(K\entails\exists x\bigwedge_i\operatorname{sgn}P_i(\bar y,x)=\sigma_i\)
iff \(\sigma\) is a letter of \(\operatorname{SIGN}(\mathcal P(\bar y))\), iff
\(K\entails\bigvee\{\theta_{\mathcal P,t}(\bar y):t\in\mathbb T_{s,D},\
\sigma\text{ is a letter of }t\}\) by
\cref{thm:p03-w03-p03-037-uniform-sign}.  The disjunction is finite (and is
\(0\neq0\) if empty) and quantifier-free, and the equivalence holds in every
model and for every \(\bar y\); so it is a consequence of \(\mathsf{RCF}\).
\end{proof}
\end{annexitem}

\begin{annexitem}{annex:p03-w03-p03-037-complete}{thm:p03-w03-p03-037-complete}{Completeness}
\begin{proof}
Let \(\varphi\) be a sentence.  By \cref{thm:p03-w03-p03-037-qe} there is a
quantifier-free sentence \(\psi\) with
\(\mathsf{RCF}\entails\varphi\leftrightarrow\psi\).  The atoms of \(\psi\)
compare closed terms, which denote \(n\cdot1\) for integers \(n\); in every
ordered field \(n\cdot1=0\) iff \(n=0\), and \(0<n\cdot1\) iff \(n>0\).  So
every atom of \(\psi\), and hence \(\psi\), has the same truth value in all
ordered fields.  \(\mathsf{RCF}\) has the model \(\mathbb R\), by the
intermediate-value theorem for continuous functions on \(\mathbb R\).  If
\(\psi\) is true in \(\mathbb R\), it is true in every model, and
\(\mathsf{RCF}\entails\varphi\); otherwise \(\mathsf{RCF}\entails\neg\varphi\).
Thus \(\mathsf{RCF}\) is complete, and every real closed field satisfies the
same sentences as \(\mathbb R\).
\end{proof}
\end{annexitem}

\begin{annexitem}{annex:p03-w03-p03-037-definable-semialgebraic}{thm:p03-w03-p03-037-definable-semialgebraic}{Definable sets are semialgebraic}
\begin{proof}
(\(\Leftarrow\)) A polynomial \(P\in K[X_1,\dots,X_n]\) with coefficients
\(c_1,\dots,c_r\) is \(\tilde P(\bar X,\bar c)\) for an integer polynomial
\(\tilde P(\bar X,\bar Z)\), obtained by replacing each coefficient by a new
variable.  So \(\{P>0\}\) and \(\{P=0\}\) are defined by the formulas
\(\tilde P(\bar x,\bar c)>0\) and \(\tilde P(\bar x,\bar c)=0\) with
parameters \(\bar c\), and Boolean combinations of sets correspond to
Boolean combinations of formulas.

(\(\Rightarrow\)) Let \(S=\{\bar x\in K^n:K\entails\varphi(\bar x,\bar c)\}\).
By \cref{thm:p03-w03-p03-037-qe} there is a quantifier-free
\(\psi(\bar x,\bar z)\) with
\(\mathsf{RCF}\entails\forall\bar x\,\forall\bar z\,(\varphi\leftrightarrow\psi)\),
and by the first step of its proof \(\psi\) is equivalent to a Boolean
combination of sign conditions \(\tilde P(\bar x,\bar z)=0\),
\(\tilde P(\bar x,\bar z)>0\) with integer polynomials \(\tilde P\).  Then
\(S\) is the same Boolean combination of the sets
\(\{\tilde P(\bar x,\bar c)=0\}\) and \(\{\tilde P(\bar x,\bar c)>0\}\), and
\(\tilde P(\bar X,\bar c)\in K[\bar X]\).  (A set \(\{F<0\}\) is
\(\{-F>0\}\).)
\end{proof}
\end{annexitem}

\begin{annexitem}{annex:p03-w03-p03-037-o-minimal}{thm:p03-w03-p03-037-o-minimal}{Every real closed field is o-minimal}
\begin{proof}
\((K,<)\) is a nonempty dense linear order without endpoints: if \(a<b\) then
\(a<(a+b)/2<b\), and \(a-1<a<a+1\).

Let \(S\subseteq K\) be defined by \(\varphi(x,\bar c)\).  By
\cref{thm:p03-w03-p03-037-qe} and the first step of its proof,
\(\varphi(x,\bar y)\) is \(\mathsf{RCF}\)-equivalent to a Boolean
combination of sign conditions on integer polynomials
\(P_1,\dots,P_s\in\mathbb Z[x,\bar y]\); let \(D\) be the sum of their
\(x\)-degrees.  Let \(\mathcal P=(P_1(x,\bar c),\dots,P_s(x,\bar c))\).  Each
\(P_i(x,\bar c)\) has constant sign on every cell of \(\mathcal P\), so each
set \(\{x:\operatorname{sgn}P_i(x,\bar c)=e\}\) is a union of cells; since the
cells partition \(K\), every Boolean combination of such sets is a union of
cells.  So \(S\) is a union of cells, which are points and open intervals with
endpoints in \(K\cup\{\pm\infty\}\), and there are finitely many.  Hence
\(K\) is o-minimal; with \cref{thm:p03-w03-p03-037-complete},
\(\mathsf{RCF}\) is a complete theory all of whose models are o-minimal.

For the bound, the number \(N\) of roots of \(\mathcal P\) is at most \(D\)
by \cref{thm:p03-w03-p03-037-poly-basics}(1).  List the cells in order as
\(C_0<C_1<\dots<C_{2N}\) (every element of \(C_i\) is below every element of
\(C_{i+1}\)).  Group the cells contained in \(S\) into maximal runs of
consecutive cells.  The union of a run \(C_i,\dots,C_{i'}\) is convex: if
\(x<z<x'\) with \(x\in C_i\) and \(x'\in C_{i'}\), then \(z\) lies in some
\(C_l\) with \(i\leq l\leq i'\).  Two different runs are separated by at least
one cell not contained in \(S\), so \(r\) runs need at least \(2r-1\) cells;
hence \(2r-1\leq2N+1\) and \(r\leq N+1\leq D+1\).  The unions of the runs are
pairwise disjoint.
\end{proof}
\end{annexitem}

\begin{annexitem}{annex:p03-w03-p03-037-unstable}{thm:p03-w03-p03-037-unstable}{Real closed fields are unstable}
\begin{proof}
Let \(K\) be any real closed field and put \(a_i=b_i=i\cdot1\) for
\(i<\omega\).  If \(i<j\) then \(b_j-a_i=(j-i)\cdot1>0\), so
\(K\entails a_i<b_j\).  If \(i\geq j\) then \(a_i-b_j=(i-j)\cdot1\geq0\), so
\(K\not\entails a_i<b_j\).  Thus \(K\entails a_i<b_j\) iff \(i<j\), which is
the order property of \cref{def:p03-stability} for \(x<y\) in the complete
theory \(\mathsf{RCF}\) (\cref{thm:p03-w03-p03-037-complete}).
\end{proof}
\end{annexitem}

\begin{annexitem}{annex:p03-w03-p03-037-nip-one-variable}{thm:p03-w03-p03-037-nip-one-variable}{One-variable formulas are NIP with an explicit bound}
\begin{proof}
Let \(K\) be a real closed field and suppose the family shatters
\(A=\{a_1<\dots<a_{2D+3}\}\subseteq K\).  Then some \(\bar b\) satisfies
\(\varphi(K;\bar b)\cap A=A_0\), where \(A_0=\{a_1,a_3,\dots,a_{2D+3}\}\)
has \(D+2\) elements.  By \cref{thm:p03-w03-p03-037-o-minimal},
\(\varphi(K;\bar b)\) is a union of at most \(D+1\) convex sets, so two
elements \(a_{2i-1}<a_{2j-1}\) of \(A_0\), with \(i<j\), lie in one convex
set \(C\subseteq\varphi(K;\bar b)\).  Since \(2i-1<2i<2j-1\),
\(a_{2i}\in C\), so \(a_{2i}\in\varphi(K;\bar b)\cap A=A_0\), which is false
because \(2i\) is even.  A family that shatters a set shatters every subset
of it, so no set with at least \(2D+3\) elements is shattered in any model.
Hence arbitrarily large finite sets are not shattered, and \(\varphi\) is
NIP in the complete theory \(\mathsf{RCF}\).
\end{proof}
\end{annexitem}

\begin{annexitem}{annex:p03-w03-p03-037-linear-shatter}{thm:p03-w03-p03-037-linear-shatter}{Positivity sets in a finite-dimensional space}
\begin{proof}
Let \(g_1,\dots,g_N\) span \(V\) and let \(A=\{x_0,\dots,x_N\}\subseteq X\)
have \(N+1\) distinct elements.  The evaluation map
\[
e:V\to F^{N+1},\qquad e(f)=(f(x_0),\dots,f(x_N)),
\]
is linear, so its image
\(W\) is spanned by \(e(g_1),\dots,e(g_N)\).  The homogeneous linear system
\(\sum_{i=0}^{N}c_i\,g_l(x_i)=0\) (\(l=1,\dots,N\)) has \(N\) equations in the
\(N+1\) unknowns \(c_0,\dots,c_N\), so Gaussian elimination over the field
\(F\) leaves at least one free unknown and gives a solution
\(c\neq0\).  By linearity \(\sum_ic_i\,w_i=0\) for every \(w\in W\), that is,
\[
\sum_{i=0}^{N}c_i\,f(x_i)=0\qquad\text{for every }f\in V.\tag{\(\dagger\)}
\]
(When \(N=0\), \(V=\{0\}\) and any \(c_0\neq0\) works.)  Replacing \(c\) by
\(-c\) if necessary, the set \(S=\{i:c_i>0\}\) is nonempty.  Let
\(A_0=\{x_i:i\in S\}\), and suppose that some \(f\in V\) satisfies
\(\{x:f(x)>0\}\cap A=A_0\).  For \(i\in S\), \(c_i>0\) and \(f(x_i)>0\), so
\(c_if(x_i)>0\).  For \(i\notin S\), \(c_i\leq0\), and \(x_i\notin A_0\)
gives \(f(x_i)\leq0\), so \(c_if(x_i)\geq0\).  Since \(S\neq\varnothing\), the
sum in \((\dagger)\) is positive, which contradicts \((\dagger)\).  So the
subset \(A_0\) of \(A\) is not cut out by any member of the family, and \(A\)
is not shattered.
\end{proof}
\end{annexitem}

\begin{annexitem}{annex:p03-w03-p03-037-sauer-shelah}{thm:p03-w03-p03-037-sauer-shelah}{Counting traces of a family of bounded shattering}
\begin{proof}
For a family \(\mathcal G\) of subsets of a finite set \(A\), let
\(\operatorname{sh}(\mathcal G)\) be the set of subsets \(B\subseteq A\)
shattered by \(\mathcal G\), that is, with
\(\{G\cap B:G\in\mathcal G\}\) equal to the set of all subsets of \(B\).

\emph{Claim.}  \(\abs{\mathcal G}\leq\abs{\operatorname{sh}(\mathcal G)}\)
for every family \(\mathcal G\) of subsets of a finite set \(A\).  By
induction on \(\abs A\).  If \(A=\varnothing\), then either
\(\mathcal G=\varnothing\), and \(0\leq\abs{\operatorname{sh}(\mathcal G)}\),
or \(\mathcal G=\{\varnothing\}\), which shatters \(\varnothing\), so both
sides are \(1\).  If \(A\neq\varnothing\), fix \(a\in A\) and put
\(A'=A\setminus\{a\}\),
\[
\mathcal G_0=\{G\in\mathcal G:a\notin G\},\qquad
\mathcal G_1=\{G\setminus\{a\}:G\in\mathcal G,\ a\in G\},
\]
both families of subsets of \(A'\).  The map \(G\mapsto G\setminus\{a\}\) is
injective on \(\{G\in\mathcal G:a\in G\}\), so
\(\abs{\mathcal G}=\abs{\mathcal G_0}+\abs{\mathcal G_1}\), and by the
induction hypothesis
\(\abs{\mathcal G}\leq\abs{\operatorname{sh}(\mathcal G_0)}+
\abs{\operatorname{sh}(\mathcal G_1)}=\abs{\mathcal U}+\abs{\mathcal I}\), where
\(\mathcal U=\operatorname{sh}(\mathcal G_0)\cup\operatorname{sh}(\mathcal G_1)\)
and
\(\mathcal I=\operatorname{sh}(\mathcal G_0)\cap\operatorname{sh}(\mathcal G_1)\).
Every \(B\in\mathcal U\) is a subset of \(A'\) shattered by \(\mathcal G\):
for \(B\subseteq A'\), \(G\cap B=(G\setminus\{a\})\cap B\), so every trace on
\(B\) of a member of \(\mathcal G_0\) or \(\mathcal G_1\) is a trace of a member
of \(\mathcal G\).  For \(B\in\mathcal I\), the set \(B\cup\{a\}\) is shattered
by \(\mathcal G\): given \(C\subseteq B\cup\{a\}\), if \(a\in C\) choose
\(H\in\mathcal G_1\) with \(H\cap B=C\setminus\{a\}\); then
\(H\cup\{a\}\in\mathcal G\) and \((H\cup\{a\})\cap(B\cup\{a\})=C\).  If
\(a\notin C\) choose \(G\in\mathcal G_0\) with \(G\cap B=C\); then
\(G\cap(B\cup\{a\})=C\) since \(a\notin G\).  The sets in \(\mathcal U\) do
not contain \(a\) and the sets \(B\cup\{a\}\) (\(B\in\mathcal I\)) do, and
\(B\mapsto B\cup\{a\}\) is injective on subsets of \(A'\).  So
\(\abs{\mathcal U}+\abs{\mathcal I}\leq\abs{\operatorname{sh}(\mathcal G)}\),
which proves the claim.

\emph{The bound.}  Let \(\mathcal G=\mathcal F\cap A\).  A subset
\(B\subseteq A\) is shattered by \(\mathcal G\) iff it is shattered by
\(\mathcal F\), because \((S\cap A)\cap B=S\cap B\) for \(B\subseteq A\).  A
family that shatters a set shatters each of its subsets, so by hypothesis
\(\mathcal F\) shatters no subset with more than \(d\) elements.  Hence
\(\operatorname{sh}(\mathcal G)\) consists of subsets of \(A\) with at most
\(d\) elements, and by the claim
\(\abs{\mathcal F\cap A}\leq\sum_{i=0}^{d}\binom ni\).  Finally
\(\binom ni=n(n-1)\cdots(n-i+1)/i!\leq n^i\), and since
\(\binom di\geq1\) for \(i\leq d\),
\(\sum_{i\leq d}n^i\leq\sum_{i\leq d}\binom di n^i=(n+1)^d\).
\end{proof}
\end{annexitem}

\begin{annexitem}{annex:p03-w03-p03-037-nip}{thm:p03-w03-p03-037-nip}{Real closed fields are NIP}
\begin{proof}
Let \(\ell=\abs{\bar x}\).  By the first step of the proof of
\cref{thm:p03-w03-p03-037-qe}, \(\psi\) is \(\mathsf{RCF}\)-equivalent to a
Boolean combination \(\beta(\alpha_1,\dots,\alpha_k)\) of atomic sign
conditions, where \(\alpha_j\) is \(P_j>0\) or \(P_j=0\); a polynomial may
occur in several atoms, so the list \(P_1,\dots,P_k\) may have repetitions.
Since \(\varphi\leftrightarrow\psi\) holds in every model, it suffices to
count traces of the family \(\{\psi(K;\bar b)\}\).

Fix a real closed field \(K\) and a finite \(A\subseteq K^\ell\) with \(n\)
elements.  For \(j\leq k\) let \(V_j\) be the \(K\)-vector space of functions
\(K^\ell\to K\) spanned by the \(N_j\) monomial functions
\(\bar a\mapsto\bar a^{\alpha}\) (\(\alpha\in I_j\)), and let
\(\mathcal G_j=\{\{\bar a:f(\bar a)>0\}:f\in V_j\}\).  By
\cref{thm:p03-w03-p03-037-linear-shatter}, with \(F=K\) and \(X=K^\ell\),
\(\mathcal G_j\) shatters no set with \(N_j+1\) elements, so by
\cref{thm:p03-w03-p03-037-sauer-shelah}
\(\abs{\mathcal G_j\cap A}\leq(n+1)^{N_j}\).

For \(\bar b\in K^{\abs{\bar y}}\) put
\(f_{j,\bar b}(\bar a)=P_j(\bar a,\bar b)=
\sum_{\alpha\in I_j}q_{j,\alpha}(\bar b)\,\bar a^{\alpha}\).  Its coefficients
\(q_{j,\alpha}(\bar b)\) lie in \(K\), so \(f_{j,\bar b}\in V_j\) and
\(-f_{j,\bar b}\in V_j\).  If \(\alpha_j\) is \(P_j>0\), then
\(\alpha_j(K;\bar b)=\{f_{j,\bar b}>0\}\in\mathcal G_j\).  If \(\alpha_j\) is
\(P_j=0\), then by trichotomy
\(\alpha_j(K;\bar b)=K^\ell\setminus(\{f_{j,\bar b}>0\}\cup\{-f_{j,\bar b}>0\})\),
so \(\alpha_j(K;\bar b)\cap A\) is determined by the pair
\((\{f_{j,\bar b}>0\}\cap A,\{-f_{j,\bar b}>0\}\cap A)\in(\mathcal G_j\cap A)^2\).
In either case, as \(\bar b\) varies, \(\alpha_j(K;\bar b)\cap A\) takes at
most \(\abs{\mathcal G_j\cap A}^2\leq(n+1)^{2N_j}\) values.

Intersection with \(A\) commutes with the Boolean operations:
\((S\cup S')\cap A=(S\cap A)\cup(S'\cap A)\),
\((S\cap S')\cap A=(S\cap A)\cap(S'\cap A)\) and
\((K^\ell\setminus S)\cap A=A\setminus(S\cap A)\).  Hence
\(\psi(K;\bar b)\cap A\) is the Boolean combination \(\beta\), taken relative
to \(A\), of the sets \(\alpha_j(K;\bar b)\cap A\), and is determined by the
tuple of these sets.  Therefore
\[
\abs{\{\psi(K;\bar b)\cap A:\bar b\in K^{\abs{\bar y}}\}}\leq
\prod_{j\leq k}(n+1)^{2N_j}=(n+1)^{d}.
\]

If the family shatters \(A\), the left side is \(2^n\), so
\(2^n\leq(n+1)^d\).  We show that \((n+1)^d<2^n\) whenever \(n\geq n_d\).
First, \(2^{d+1}\geq2d+2\) for every \(d\geq0\) (true for \(d=0\), and
\(2^{d+2}\geq4d+4\geq2(d+1)+2\)), so \(n\geq n_d\geq2^{d+1}\geq2d+2\).  Then
\(n\geq d+1\), and
\[
2^n=\sum_{i=0}^{n}\binom ni\geq\binom n{d+1}
=\frac{n(n-1)\cdots(n-d)}{(d+1)!}\geq\frac{(n-d)^{d+1}}{(d+1)!}.
\]
Since \(n\geq2d+1\), \(n-d\geq(n+1)/2\), so
\[
2^n\geq\frac{(n+1)^{d+1}}{2^{d+1}(d+1)!}=(n+1)^d\cdot\frac{n+1}{n_d}>(n+1)^d,
\]
because \(n+1>n\geq n_d\).  So no set with \(n\geq n_d\) elements is
shattered in any real closed field.  In particular arbitrarily large finite
sets are not shattered in any model of \(\mathsf{RCF}\), so \(\varphi\) is NIP
in the complete theory \(\mathsf{RCF}\)
(\cref{thm:p03-w03-p03-037-complete}); as \(\varphi\) was arbitrary,
\(\mathsf{RCF}\) is NIP.
\end{proof}
\end{annexitem}

\begin{annexitem}{annex:p03-w03-p03-037-ring-language}{thm:p03-w03-p03-037-ring-language}{Elimination needs the order symbol}
\begin{proof}
A quantifier-free formula \(\psi(x)\) of \(\{0,1,+,-,\mathbin{\cdot}\}\)
without parameters is a Boolean combination of equations \(t_1=t_2\), each
equivalent in a ring to \(P(x)=0\) for an integer polynomial \(P\).  In
\(K\), which has characteristic \(0\), \(P\) is the zero polynomial or has
finitely many roots (\cref{thm:p03-w03-p03-037-poly-basics}(1)); so each
atom defines \(K\) or a finite set, and every Boolean combination defines a
finite or a cofinite subset of \(K\).

The formula \(\exists z\,(x=z\mathbin{\cdot}z)\) defines the set of squares,
which is \(\{x:x\geq0\}\): squares are nonnegative, and every nonnegative
element is a square by \cref{thm:p03-w03-p03-037-poly-basics}(7).  This set
contains the distinct elements \(n\cdot1\) (\(n\geq0\)) and its complement
contains \(-n\cdot1\) (\(n\geq1\)), so it is neither finite nor cofinite.
Hence no quantifier-free \(\psi(x)\) is equivalent to it in \(K\), and a
fortiori no such equivalence follows from the theory of the ring reduct.
\end{proof}
\end{annexitem}

\begin{annexitem}{annex:p03-w03-p03-037-expansion}{thm:p03-w03-p03-037-expansion}{An expansion that is neither o-minimal nor NIP}
\begin{proof}
\emph{Not o-minimal.}  The structure expands the dense linear order
\((\mathbb R,<)\) without endpoints, and \(Z(x)\) defines \(\mathbb Z\).
Suppose \(\mathbb Z\) were a finite union of points and open intervals.  No
open interval \((\alpha,\beta)\) is contained in \(\mathbb Z\): pick
\(a\in(\alpha,\beta)\) and \(h\) with \(0<h<1\) and
\((a-h,a+h)\subseteq(\alpha,\beta)\); if \(a\notin\mathbb Z\) we are done,
and otherwise \(a+h/2\in(a,a+1)\) is not an integer.  So the union would
consist of finitely many points, while \(\mathbb Z\) is infinite.

\emph{Independence property.}  The complete theory is
\(\operatorname{Th}(\mathcal R_{\mathbb Z})\), and \(\mathcal R_{\mathbb Z}\)
is one of its models.  For a positive integer \(n\) let \(p_1<\dots<p_n\) be
primes, which exist because there are infinitely many primes, and let
\(A=\{p_1,\dots,p_n\}\).  For \(I\subseteq A\) let \(b_I=\prod_{q\in I}q\),
with \(b_\varnothing=1\).  Then \(\mathcal R_{\mathbb Z}\entails\delta(p;b_I)\)
iff \(p\) divides \(b_I\) in \(\mathbb Z\).  If \(p\in I\), it does.  If
\(p\notin I\) and \(p\mid b_I\), then \(I\neq\varnothing\) because \(p\geq2\)
does not divide \(1\), and Euclid's lemma (a prime dividing a product of
integers divides one of the factors) gives \(p\mid q\) for some prime
\(q\in I\), so \(p=q\in I\), a contradiction.  Hence
\(\delta(\mathcal R_{\mathbb Z};b_I)\cap A=I\) for every \(I\subseteq A\):
\(A\) is shattered.  As \(n\) is arbitrary, \(\delta\) has the independence
property in the sense of \cref{def:p03-nip}.
\end{proof}
\end{annexitem}
